\section{AMC 8 1994}
        \begin{enumerate}
        	\item (\textit{AMC 8 1994}) Which of the following is the largest?


 $\text{(A)}\ \dfrac{1}{3} \qquad \text{(B)}\ \dfrac{1}{4} \qquad \text{(C)}\ \dfrac{3}{8} \qquad \text{(D)}\ \dfrac{5}{12} \qquad \text{(E)}\ \dfrac{7}{24}$


 

	\item (\textit{AMC 8 1994}) $\dfrac{1}{10}+\dfrac{2}{10}+\dfrac{3}{10}+\dfrac{4}{10}+\dfrac{5}{10}+\dfrac{6}{10}+\dfrac{7}{10}+\dfrac{8}{10}+\dfrac{9}{10}+\dfrac{55}{10}=$


 $\text{(A)}\ 4\dfrac{1}{2} \qquad \text{(B)}\ 6.4 \qquad \text{(C)}\ 9 \qquad \text{(D)}\ 10 \qquad \text{(E)}\ 11$


 

	\item (\textit{AMC 8 1994}) Each day Maria must work $8$ hours.  This does not include the $45$ minutes she takes for lunch.  If she begins working at $\text{7:25 A.M.}$ and takes her lunch break at noon, then her working day will end at


 $\text{(A)}\ \text{3:40 P.M.} \qquad \text{(B)}\ \text{3:55 P.M.} \qquad \text{(C)}\ \text{4:10 P.M.} \qquad \text{(D)}\ \text{4:25 P.M.} \qquad \text{(E)}\ \text{4:40 P.M.}$


 

	\item (\textit{AMC 8 1994}) Which of the following represents the result when the figure shown below is rotated clockwise $120^\circ$ around its center?


 

                    \begin{figure}[H]
                        \centering
                        \includegraphics[width=0.5\textwidth]{Problem/AMC_8_AJHSME_1994_4_Problem_Figure_1.png}
                    \end{figure}
                    


 


 
                    \begin{figure}[H]
                        \centering
                        \includegraphics[width=0.5\textwidth]{Problem/AMC_8_AJHSME_1994_4_Problem_Figure_2.png}
                    \end{figure}
                    


 

	\item (\textit{AMC 8 1994}) Given that $\text{1 mile} = \text{8 furlongs}$ and $\text{1 furlong} = \text{40 rods}$, the number of rods in one mile is


 $\text{(A)}\ 5 \qquad \text{(B)}\ 320 \qquad \text{(C)}\ 660 \qquad \text{(D)}\ 1760 \qquad \text{(E)}\ 5280$


 

	\item (\textit{AMC 8 1994}) The unit's digit (one's digit) of the product of any six consecutive positive whole numbers is


 $\text{(A)}\ 0 \qquad \text{(B)}\ 2 \qquad \text{(C)}\ 4 \qquad \text{(D)}\ 6 \qquad \text{(E)}\ 8$


 

	\item (\textit{AMC 8 1994}) If $\angle A = 60^\circ$, $\angle E = 40^\circ$ and $\angle C = 30^\circ$, then $\angle BDC =$


 

                    \begin{figure}[H]
                        \centering
                        \includegraphics[width=0.5\textwidth]{Problem/AMC_8_AJHSME_1994_7_Problem_Figure_1.png}
                    \end{figure}
                    


 


 $\text{(A)}\ 40^\circ \qquad \text{(B)}\ 50^\circ \qquad \text{(C)}\ 60^\circ \qquad \text{(D)}\ 70^\circ \qquad \text{(E)}\ 80^\circ$


 

	\item (\textit{AMC 8 1994}) For how many three-digit whole numbers does the sum of the digits equal $25$?


 $\text{(A)}\ 2 \qquad \text{(B)}\ 4 \qquad \text{(C)}\ 6 \qquad \text{(D)}\ 8 \qquad \text{(E)}\ 10$


 

	\item (\textit{AMC 8 1994}) A shopper buys a $100$ dollar coat on sale for $20\%$ off.  An additional $5$ dollars are taken off the sale price by using a discount coupon.  A sales tax of $8\%$ is paid on the final selling price.  The total amount the shopper pays for the coat is


 $\text{(A)}\ \text{81.00 dollars} \qquad \text{(B)}\ \text{81.40 dollars} \qquad \text{(C)}\ \text{82.00 dollars} \qquad \text{(D)}\ \text{82.08 dollars} \qquad \text{(E)}\ \text{82.40 dollars}$


 

	\item (\textit{AMC 8 1994}) For how many positive integer values of $N$ is the expression $\dfrac{36}{N+2}$ an integer?


 $\text{(A)}\ 7 \qquad \text{(B)}\ 8 \qquad \text{(C)}\ 9 \qquad \text{(D)}\ 10 \qquad \text{(E)}\ 12$


 

	\item (\textit{AMC 8 1994}) Last summer $100$ students attended basketball camp.  Of those attending, $52$ were boys and $48$ were girls.  Also, $40$ students were from Jonas Middle School and $60$ were from Clay Middle School.  Twenty of the girls were from Jonas Middle School.  How many of the boys were from Clay Middle School?


 $\text{(A)}\ 20 \qquad \text{(B)}\ 32 \qquad \text{(C)}\ 40 \qquad \text{(D)}\ 48 \qquad \text{(E)}\ 52$


 

	\item (\textit{AMC 8 1994}) Each of the three large squares shown below is the same size.  Segments that intersect the sides of the squares intersect at the midpoints of the sides.  How do the shaded areas of these squares compare?


 

                    \begin{figure}[H]
                        \centering
                        \includegraphics[width=0.5\textwidth]{Problem/AMC_8_AJHSME_1994_12_Problem_Figure_1.png}
                    \end{figure}
                    


 


 $\text{(A)}\ \text{The shaded areas in all three are equal.}$


 $\text{(B)}\ \text{Only the shaded areas of }I\text{ and }II\text{ are equal.}$


 $\text{(C)}\ \text{Only the shaded areas of }I\text{ and }III\text{ are equal.}$


 $\text{(D)}\ \text{Only the shaded areas of }II\text{ and }III\text{ are equal.}$


 $\text{(E)}\ \text{The shaded areas of }I, II\text{ and }III\text{ are all different.}$


 

	\item (\textit{AMC 8 1994}) The number halfway between $\dfrac{1}{6}$ and $\dfrac{1}{4}$ is


 $\text{(A)}\ \dfrac{1}{10} \qquad \text{(B)}\ \dfrac{1}{5} \qquad \text{(C)}\ \dfrac{5}{24} \qquad \text{(D)}\ \dfrac{7}{24} \qquad \text{(E)}\ \dfrac{5}{12}$


 

	\item (\textit{AMC 8 1994}) Two children at a time can play pairball.  For $90$ minutes, with only two children playing at a time, five children take turns so that each one plays the same amount of time.  The number of minutes each child plays is


 $\text{(A)}\ 9 \qquad \text{(B)}\ 10 \qquad \text{(C)}\ 18 \qquad \text{(D)}\ 20 \qquad \text{(E)}\ 36$


 

	\item (\textit{AMC 8 1994}) If this path is to continue in the same pattern:


 

                    \begin{figure}[H]
                        \centering
                        \includegraphics[width=0.5\textwidth]{Problem/AMC_8_AJHSME_1994_15_Problem_Figure_1.png}
                    \end{figure}
                    


 


 then which sequence of arrows goes from point $425$ to point $427$?


 
                    \begin{figure}[H]
                        \centering
                        \includegraphics[width=0.5\textwidth]{Problem/AMC_8_AJHSME_1994_15_Problem_Figure_2.png}
                    \end{figure}
                    


 

	\item (\textit{AMC 8 1994}) The perimeter of one square is $3$ times the perimeter of another square.  The area of the larger square is how many times the area of the smaller square?


 $\text{(A)}\ 2 \qquad \text{(B)}\ 3 \qquad \text{(C)}\ 4 \qquad \text{(D)}\ 6 \qquad \text{(E)}\ 9$


 

	\item (\textit{AMC 8 1994}) Pauline Bunyan can shovel snow at the rate of $20$ cubic yards for the first hour, $19$ cubic yards for the second, $18$ for the third, etc., always shoveling one cubic yard less per hour than the previous hour.  If her driveway is $4$ yards wide, $10$ yards long, and covered with snow $3$ yards deep, then the number of hours it will take her to shovel it clean is closest to


 $\text{(A)}\ 4 \qquad \text{(B)}\ 5 \qquad \text{(C)}\ 6 \qquad \text{(D)}\ 7 \qquad \text{(E)}\ 12$


 

	\item (\textit{AMC 8 1994}) Mike leaves home and drives slowly east through city traffic.  When he reaches the highway he drives east more rapidly until he reaches the shopping mall where he stops.  He shops at the mall for an hour.  Mike returns home by the same route as he came, driving west rapidly along the highway and then slowly through city traffic.  Each graph shows the distance from home on the vertical axis versus the time elapsed since leaving home on the horizontal axis.  Which graph is the best representation of Mike's trip?


 

                    \begin{figure}[H]
                        \centering
                        \includegraphics[width=0.5\textwidth]{Problem/AMC_8_AJHSME_1994_18_Problem_Figure_1.png}
                    \end{figure}
                    


 


 

	\item (\textit{AMC 8 1994}) Around the outside of a $4$ by $4$ square, construct four semicircles (as shown in the figure) with the four sides of the square as their diameters.  Another square, $ABCD$, has its sides parallel to the corresponding sides of the original square, and each side of $ABCD$ is tangent to one of the semicircles.  The area of the square $ABCD$ is


 

                    \begin{figure}[H]
                        \centering
                        \includegraphics[width=0.5\textwidth]{Problem/AMC_8_AJHSME_1994_19_Problem_Figure_1.png}
                    \end{figure}
                    


 


 $\text{(A)}\ 16 \qquad \text{(B)}\ 32 \qquad \text{(C)}\ 36 \qquad \text{(D)}\ 48 \qquad \text{(E)}\ 64$


 

	\item (\textit{AMC 8 1994}) Let $W,X,Y$ and $Z$ be four different digits selected from the set


 $\{ 1,2,3,4,5,6,7,8,9\}.$

 If the sum $\dfrac{W}{X} + \dfrac{Y}{Z}$ is to be as small as possible, then $\dfrac{W}{X} + \dfrac{Y}{Z}$ must equal


 $\text{(A)}\ \dfrac{2}{17} \qquad \text{(B)}\ \dfrac{3}{17} \qquad \text{(C)}\ \dfrac{17}{72} \qquad \text{(D)}\ \dfrac{25}{72} \qquad \text{(E)}\ \dfrac{13}{36}$


 

	\item (\textit{AMC 8 1994}) A gumball machine contains $9$ red, $7$ white, and $8$ blue gumballs.  The least number of gumballs a person must buy to be sure of getting four gumballs of the same color is


 $\text{(A)}\ 8 \qquad \text{(B)}\ 9 \qquad \text{(C)}\ 10 \qquad \text{(D)}\ 12 \qquad \text{(E)}\ 18$


 

	\item (\textit{AMC 8 1994}) The two wheels shown below are spun and the two resulting numbers are added.  The probability that the sum is even is


 

                    \begin{figure}[H]
                        \centering
                        \includegraphics[width=0.5\textwidth]{Problem/AMC_8_AJHSME_1994_22_Problem_Figure_1.png}
                    \end{figure}
                    


 


 $\text{(A)}\ \dfrac{1}{6} \qquad \text{(B)}\ \dfrac{1}{4} \qquad \text{(C)}\ \dfrac{1}{3} \qquad \text{(D)}\ \dfrac{5}{12} \qquad \text{(E)}\ \dfrac{4}{9}$


 

	\item (\textit{AMC 8 1994}) If $X$, $Y$ and $Z$ are different digits, then the largest possible $3-$digit sum for


 
$\begin{tabular}{ccc} X & X & X \\  & Y & X \\ + & & X \\ \hline \end{tabular}$


 


 has the form


 $\text{(A)}\ XXY \qquad \text{(B)}\ XYZ \qquad \text{(C)}\ YYX \qquad \text{(D)}\ YYZ \qquad \text{(E)}\ ZZY$


 

	\item (\textit{AMC 8 1994}) A $2$ by $2$ square is divided into four $1$ by $1$ squares.  Each of the small squares is to be painted either green or red.  In how many different ways can the painting be accomplished so that no green square shares its top or right side with any red square?  There may be as few as zero or as many as four small green squares.


 $\text{(A)}\  4 \qquad \text{(B)}\ 6 \qquad \text{(C)}\ 7 \qquad \text{(D)}\ 8 \qquad \text{(E)}\ 16$


 

	\item (\textit{AMC 8 1994}) Find the sum of the digits in the answer to


 
$\underbrace{9999\ldots 99}_{94\text{ nines}} \times \underbrace{4444\ldots 44}_{94\text{ fours}}$


 


 where a string of $94$ nines is multiplied by a string of $94$ fours.


 $\text{(A)}\ 846 \qquad \text{(B)}\ 855 \qquad \text{(C)}\ 945 \qquad \text{(D)}\ 954 \qquad \text{(E)}\ 1072$


 

\end{enumerate}